\section{Periods containing a rational direction}
\label{rat:section}

Let
\[
 \Lambda=3U\oplus 2E_8(-1),\qquad
 \mathcal D=\{P\subset\Lambda\otimes\R:
       P\text{ is an oriented positive two-plane}\}.
\]
Here positivity refers to the intersection form, denoted by $(\ ,\ )$.
For a primitive vector $v\in\Lambda$ with $(v,v)>0$, set
\[
 \mathcal D_v=\{P\in\mathcal D:v\in P\}.
\]
All open subsets of $\mathcal D_v$ below have its relative real topology.
We use the K3 period, Torelli, and lattice theorems in their usual
unpolarized form; a reference for these results and the facts about
genus-one pencils used below is
\cite[Chapters~2, 6--8, and~11]{Huybrechts}.

\begin{theorem}\label{rat:open}
For every primitive positive integral vector $v\in\Lambda$, there is a
nonempty relatively open subset $\mathcal O_v\subset\mathcal D_v$ such
that every marked K3 surface with period in $\mathcal O_v$ has property
$L$ at every point.
\end{theorem}

On a marked K3 surface whose period contains $v$, the holomorphic
symplectic form has a unique normalization satisfying
\begin{equation}\label{rat:normalization}
 [\operatorname{Im}\sigma]=v.
\end{equation}
Indeed the real and imaginary parts of a nonzero holomorphic two-form
are orthogonal and of equal positive square, and multiplication by a
nonzero complex scalar gives all oriented orthogonal frames of its
period plane. This normalization varies smoothly on $\mathcal D_v$.

We first produce a convenient central surface, then construct strips
on its nearby deformations satisfying \eqref{rat:normalization}.

\subsection{An isotrivial center for every positive integral vector}

\begin{lemma}\label{rat:center}
For every primitive $v\in\Lambda$ of square $2d>0$, there is a marked
projective K3 surface $X_0$ whose period contains $v$ and which admits
a genus-one fibration
\[
 \pi:X_0\longrightarrow\mathbb P^1
\]
with all smooth fibers of $j$-invariant zero. Its primitive fiber class
$e\in\Lambda$ satisfies $(e,v)=0$.
\end{lemma}

\begin{proof}
We begin with an auxiliary elliptic K3 surface with section, given by
the Weierstrass equation
\[
 y^2=x^3+a(t),
\]
where $a$ is a section of $\mathcal O_{\mathbb P^1}(12)$ with twelve
simple zeros, and the fundamental line bundle is
$\mathcal O_{\mathbb P^1}(2)$. The canonical bundle formula gives
trivial canonical bundle, and
$R^1\pi_*\mathcal O=\mathcal O_{\mathbb P^1}(-2)$ gives irregularity
zero. The simple zeros give cuspidal fibers with smooth total space.
Thus the smooth Weierstrass surface is a K3 surface. These global
Weierstrass and direct-image formulas are recalled in
\cite[Chapter~11, Sections~1--2]{Huybrechts}.

Let $g$ be the automorphism which multiplies $x$ by a primitive cube
root of unity. Its fixed locus is the disjoint union of the zero
section and the double cover $x=0$ of the base. The latter is branched
at twelve points and has genus five. The fixed locus has Euler
characteristic $2+(2-10)=-6$. The topological Lefschetz formula, in the order-three K3 formulation
of \cite[proof of Theorem~2.2]{ArtebaniSarti}, gives
\[
 \operatorname{tr}(g\mid H^2)=-8.
\]
If $r$ is the invariant rank, the other eigenvalues occur in conjugate
pairs of primitive cube roots, so
$r-(22-r)/2=-8$, whence $r=2$. The fiber and zero section generate a
unimodular hyperbolic plane in this invariant lattice; consequently
the invariant lattice is exactly $U$. Its orthogonal complement is
even unimodular of signature $(2,18)$ and contains two hyperbolic
planes.

That complement contains a primitive vector of square $2d$.
Eichler's criterion, applied to the unimodular K3 lattice with its
two hyperbolic summands, says that primitive vectors with the same
square are isometric \cite{Eichler}. We may therefore transport the
preceding lattice isometry so that the specified vector $v$ belongs
to its noninvariant complement. There
$1+g+g^2=0$, and thus
\[
 (v,v)=(gv,gv)=2d,\qquad (v,gv)=-d.
\]
The plane $P_0=\operatorname{span}_{\R}\{v,gv\}$ is positive. Choose
either of its two orientations, and realize it as a marked K3 period
by period surjectivity. The invariant $U$ is algebraic at this period
and contains a positive integral class, so the resulting surface is
projective.

Choose a primitive isotropic generator $e$ of this $U$, with sign
chosen so that its represented class is effective. We keep the abstract
lattice data $g,U,e$ fixed and change the marking by reflections in
algebraic $(-2)$-classes so that the class represented by $e$ becomes
nef. The reflections fix $P_0$ pointwise, hence fix $v$; all preceding
lattice relations are unchanged.
For completeness, if an effective isotropic class has negative
intersection with an irreducible curve, that curve is a $(-2)$-curve;
reflection in it reduces the positive integral degree against a
fixed ample class. Repetition terminates in the nef cone. The
primitive nef isotropic pencil theorem then gives a base-point-free
complete pencil $|e|$ with
\[
 h^0(\mathcal O(e))=2,\qquad H^1(\mathcal O(e))=0.
\]
It defines a genus-one fibration on our center.

We must show that this particular pencil is isotrivial. The
$g$-invariant period planes of the chosen eigenvalue form a complex
ball: they are the positive lines in one of the two nonreal
eigenspaces of $g$ on $U^\perp\otimes\C$. Take a small ball through
$P_0$, realized by a marked local K3 family. At a generic point its
algebraic lattice is precisely $U$. To see this, each integral class
outside $U$ imposes a proper linear Hodge condition; such a condition
cannot vanish on the entire eigenspace because the two conjugate
eigenspaces span $U^\perp\otimes\C$. Their countable union has empty
interior.

The line bundle with class $e$ extends to this small family. One way
to verify the extension is to take a contractible Stein base. The
topological class extends to the total space and has zero image in
$H^2(\mathcal O)$: by the Leray spectral sequence and Stein
vanishing this group is the space of sections of
$R^2\pi_*\mathcal O$, and the class vanishes fiberwise because it is
of type $(1,1)$. The exponential sequence supplies the line bundle.
The vanishing of $H^1(\mathcal O(e))$ and cohomology and base change
extend its two sections. Their having no common zero is open by
properness. Thus the genus-one pencils persist near $P_0$.

At every nearby generic period, $g$ acts identically on the Picard
lattice, hence fixes an ample class. Global Torelli realizes it as
an automorphism of order three. Suppose its induced action on the
pencil base were nontrivial. All its fixed points would then lie
over the two fixed points of that base action. Every fixed curve
would be vertical. Since an irreducible vertical curve $C$ satisfies
$(C,e)=0$, the Hodge index theorem gives $C^2\leq0$ and adjunction
gives $p_a(C)\leq1$. The fixed locus of a finite-order automorphism
of a smooth surface is a disjoint union of smooth curves and
isolated points. It would therefore have nonnegative Euler
characteristic, contrary to the value $-6$ determined by the
cohomological action.

The base action is consequently trivial. The multiplier on a
smooth fiber's holomorphic differential equals the primitive
cube-root multiplier on the surface form. A smooth elliptic curve
with such an automorphism has $j=0$. Finally fix any smooth member
of the central pencil. It remains smooth in the extended pencil
over a neighborhood in the deformation base, and its $j$-invariant
varies continuously there. The nearby generic values are zero, so
the central value is zero as well. This proves the assertion for
every smooth member at $P_0$. The relation $(e,v)=0$ holds because
$e$ belongs to the original invariant lattice.
\end{proof}

\subsection{Actions on the smooth fibers}

Fix the surface of Lemma~\ref{rat:center} and the normalized form
$\sigma_0$. Write $S\subset\mathbb P^1$ for the finite set of
singular values. The linear monodromy of a constant-modulus family
of elliptic curves is finite. Choose a connected finite cover of
$\mathbb P^1\setminus S$ on which that monodromy is trivial, and
compactify it to a smooth compact curve $B$. Denote the resulting
punctured curve by $B^\circ$. A fixed oriented basis
$\gamma_1,\gamma_2$ now identifies the homology of all smooth
fibers over $B^\circ$.

The base change need not possess a global section. Locally it has
holomorphic sections and product torus charts. For any primitive
cycle $\gamma$ there are cylinder coordinates, with the chosen
cycle represented by the period $1$, in which
\begin{equation}\label{rat:cylinder-form}
 \sigma_0=dw\wedge dp,\qquad dp=\eta_\gamma.
\end{equation}
Here $\eta_\gamma$ is a nowhere-zero holomorphic one-form on the
corresponding base patch. Indeed, in a local product chart the
coefficient of the surface form is constant along each compact
elliptic fiber. The period normalization makes these base forms
agree on overlaps. They therefore define a global one-form
$\eta_\gamma$ on $B^\circ$. For the two basis cycles the forms
are in a constant nonreal ratio.

If $l$ is an oriented smooth loop in a K3 surface with normalized
form $\sigma$, define its action by
\begin{equation}\label{rat:action}
 H_\sigma(l)=\int_C\operatorname{Im}\sigma\pmod{\Z},
 \qquad \partial C=l.
\end{equation}
The chain $C$ exists because $H_1(X,\Z)=0$, and the value is
independent of $C$ because $[\operatorname{Im}\sigma]=v$ is
integral. On $X_0$, representatives of a fixed homology class in a
smooth elliptic fiber give the same action, since the restriction
of $\sigma_0$ to that fiber is zero. Let
\[
 h_\gamma:B^\circ\longrightarrow\R/\Z
\]
be this action function. Orientations in
\eqref{rat:cylinder-form} may be chosen so that
$dh_\gamma=\operatorname{Im}\eta_\gamma$; changing all these signs
would make no difference below. Actions are additive in the cycle.

\begin{lemma}\label{rat:safe}
The forms $\eta_\gamma$ and functions $h_\gamma$ extend to $B$.
The map
\[
 h=(h_{\gamma_1},h_{\gamma_2}):B\longrightarrow(\R/\Z)^2
\]
has finite fibers. For each primitive $\gamma$, put
\[
 F_\gamma=h_\gamma(B\setminus B^\circ).
\]
Outside a finite subset of $B$, at least two of the three cycles
$\gamma_1,\gamma_2,\gamma_1+\gamma_2$ have action outside their
respective finite sets $F_\gamma$. Any two of these cycles form an
integral homology basis.
\end{lemma}

\begin{proof}
Integration over the elliptic fibers in local product charts
expresses the pullback symplectic volume as a fixed positive
constant times $i\eta_\gamma\wedge\overline{\eta_\gamma}$.
The base change has finite degree and $X_0$ is compact, so this
one-form has finite $L^2$ norm on $B^\circ$. In a punctured local
coordinate, a holomorphic one-form with finite $L^2$ norm has no
negative Laurent coefficients. It thus extends holomorphically
across each puncture. A local primitive of the extended form
also extends the circle-valued action function.

The holomorphic forms $\eta_{\gamma_1}$ and
$\eta_{\gamma_2}$ have constant nonreal ratio. Consequently an
invertible real linear map identifies the pair of their imaginary
parts with a complex-linear differential. Endow $(\R/\Z)^2$ with this
constant complex structure. Then $h$ is a nonconstant holomorphic
map from the compact Riemann surface $B$ to an elliptic curve.
It has finite fibers.

A point without two safe choices satisfies two forbidden-level
conditions. Each pair of the integral linear functions
$h_{\gamma_1},h_{\gamma_2},h_{\gamma_1}+h_{\gamma_2}$ is an
automorphism of the real two-torus. Thus any two such conditions
restrict $h$ to a finite set of torus values. The inverse image is
finite by the preceding paragraph.
\end{proof}

We call a cycle safe at a point if its action avoids $F_\gamma$.
Every closed set of safe action values has compact inverse image
in $B^\circ$. This elementary compactness statement, rather than
a bound on the number of cylinder blocks, will control all the
successive center adjustments.

\subsection{Uniform strips after deformation}

Take a marked local deformation of $X_0$, identify its underlying
smooth fibers, and restrict its parameters to $\mathcal D_v$.
Let $X_t$ and $\sigma_t$ denote the fibers and normalized forms.
The genus-one fibration is used only on $X_0$; it is not asserted
to survive on $X_t$.

\begin{lemma}\label{rat:strips}
Fix compact families of regular torus charts on $X_0$, each
equipped with a safe primitive cycle, and suppose their actions
are separated from the corresponding forbidden values by a
positive margin. Prescribe any finite number of period cells in
these initial data. For all sufficiently small $t\in\mathcal D_v$,
each datum gives a symplectic immersed strip on $X_t$, periodic in
its long coordinate with period $1$, of positive transverse width.
On the prescribed initial cells its chart is arbitrarily close in
$C^1$ to the original cylinder chart. All widths and closeness
requirements can be chosen uniformly over the given compact data.
\end{lemma}

\begin{proof}
The centers may move after each exactness correction. We therefore
prepare the compact region in which the construction will continue
before choosing the charts. Let $\rho>0$ be a common lower bound
for the distance of the specified initial actions from their
forbidden sets. For every cycle used in the initial data, put
\[
 K_\gamma=\{b\in B:
   \operatorname{dist}(h_\gamma(b),F_\gamma)\geq\rho/4\}.
\]
This set is compact in $B^\circ$. Choose finitely many
product-torus charts covering it, with larger buffered domains
still lying in $B^\circ$. All subsequent restrictions and
recentering will use this finite atlas. We shall prove that the
updated centers remain a fixed distance inside the prepared region.

In a chart for a chosen cycle, the torus has constant lattice
$\Z+\tau\Z$, where $\operatorname{Im}\tau>0$ after choosing a
positive complementary cycle. Unroll this complementary period
while retaining $w$ modulo $\Z$. We obtain finite cylinders of
any prescribed fixed length. Starting sections can be chosen at
any fiber phase, so the chart families cover all phases, not just
a particular section.

At matching nominal base positions, successive model cylinders
are identified on their overlaps by
\begin{equation}\label{rat:model-transition}
 (w,p)\longmapsto(w+a(p),p).
\end{equation}
The functions $a$ incorporate both changes of local section and
any required number of complementary periods. The latter are
constant, because the modulus is constant. On smaller members of
the finite buffered atlas, the derivatives of the remaining
translation functions have uniform bounds. Choose long block
cores containing the prescribed cells and leave several period
cells as buffers at each end. At a joint, integral complementary
translations let the two cut heights agree at the center with a
bounded adjustment inside these buffers. Both cuts of each raw
chart may be chosen independently. After translating the long
origins, these charts have the block and overlap geometry required
by Theorem~\ref{sew:chain}. Shrinking the transverse radius makes
the variations of their cut heights uniformly smaller than the
overlap margins.

Each cylinder times its transverse disc is Stein. The local
deformation argument used earlier therefore deforms these maps
on their finite buffered domains to the nearby fibers. Although
a cylinder map is usually not injective, it is an immersion, and
on each overlap its local inverse branch close to
\eqref{rat:model-transition} is unique. The branches consequently
agree around the cylinder, including its period identification.
The resulting transition is a degree-one periodic lift. The
pulled-back forms can be normalized to $dw\wedge dp$ by the
transverse integration used in Theorem~\ref{sew:chain}.
All these constructions are uniform on the finite atlas, and
their errors tend to zero with $t$.

We next compare the action of an individual deformed block with
that of its model, before any blocks have been joined. Let $b$ be
the nominal base position of a freshly centered block, and let
$l_b$ be the image of a period loop on its zero slice. Uniformly
over the prepared charts and their allowed recenterings,
\begin{equation}\label{rat:action-comparison}
 \operatorname{dist}_{\R/\Z}
    \bigl(h_\gamma(b),H_{\sigma_t}(l_b)\bigr)<\varepsilon_t,
 \qquad \varepsilon_t\longrightarrow0.
\end{equation}
To verify this, under the fixed smooth identification the real forms
$\operatorname{Im}\sigma_t-\operatorname{Im}\sigma_0$ are exact
and tend to zero smoothly. Fix a Riemannian metric on the underlying compact manifold and let
$G$ be its Hodge Green operator. For the exact two-form
$\eta_t=\operatorname{Im}\sigma_t-\operatorname{Im}\sigma_0$, put
$\beta_t=d^*G\eta_t$. Its harmonic projection vanishes and
$dG\eta_t=Gd\eta_t=0$, so Hodge decomposition gives $d\beta_t=\eta_t$.
Elliptic regularity makes $\beta_t$ tend to zero smoothly;
see \cite[Theorem~3.2 and Equation~(3.26)]{ChenLiHodge} for this
Hodge decomposition and Green-operator formulation. Integration of
these primitives around the model loops, together with the
small cylinders between nearby model and actual loops, proves
\eqref{rat:action-comparison}. The loop families are compact:
the base positions and fiber phases range over compact sets, and
the number of wraps in a block is fixed. The estimate depends on
the comparison of this block with its own model, not on any
previous choices of centers.

Starting from any chosen block, extend the chain in both
directions. At a joint initially use a chart with matching nominal
base position. Its actual symplectic transition is close to
\eqref{rat:model-transition}. Its period defect is the constant
\[
 \int_{\R/\Z}(P\,dW-p\,dw).
\]
Moving this candidate's transverse center by the appropriately
signed defect makes the period exactly zero. Retain this
translated actual chart as the next block, so the exact overlap
just obtained is preserved. Move its nominal center by the same
amount in the model $p$ coordinate; that updated position is used
only when selecting the following candidate block. The adjustment
is small, and comparison with the model of the newly centered
chart retains its original error bound. It remains to justify
that the nominal positions never approach a puncture.

For this purpose use the action \eqref{rat:action} on the actual
surface. Its value on the period loop of the zero slice is
constant within a block. It is also constant between adjacent
blocks. Indeed the primitive of $dw\wedge dp$ is $-p\,dw$;
the change of action between the image of a zero-slice loop and
the next zero-slice loop is the imaginary part, up to sign, of
the displayed period defect. The image loop winds once in the
next cylinder. Vanishing of the defect proves equality of its
action with that of the next zero slice. Denote this common
circle value by $c$.

Thus the single-block estimate \eqref{rat:action-comparison} reads
$\operatorname{dist}_{\R/\Z}(h_\gamma(b),c)<\varepsilon_t$
at every stage. We can now verify that the next center stays in
the prepared region. Require the action error and the change in
nominal action caused by a single center adjustment to be less
than $\rho/8$. For an initial block the conserved value $c$ has
distance greater than $7\rho/8$ from $F_\gamma$. Inductively,
a next nominal center may first be chosen at the preceding one;
the small adjustment keeps it inside the prepared region. The
exact action comparison then improves its distance from the
forbidden set to greater than $3\rho/4$. This restores a fixed
margin for the following step. The same argument applies
backwards. Thus the entire doubly infinite chain is available
with common transverse and overlap margins.

We have obtained exact symplectic transitions with uniformly
small errors and the block geometry required by
Theorem~\ref{sew:chain}. That theorem supplies the periodic strip
and arbitrarily small corrections on every prescribed compact
part of the initial block. Estimates on slightly larger domains
give the asserted $C^1$ closeness. The transverse width has a
uniform positive lower bound after the fixed margin losses.
\end{proof}

\subsection{Two cycles cannot give the same leaf}

\begin{lemma}\label{rat:two-periods}
In the setting of Lemma~\ref{rat:strips}, suppose two initial
strips are obtained from independent primitive cycles of one
central smooth torus. One can choose a finite required initial block
length and a positive $C^1$ error tolerance from the central torus
and its two cycles, before restricting the deformation parameter,
with the following property. For strips meeting these bounds, if
\begin{equation}\label{rat:destroy-fiber}
 ([\sigma_t],e)\ne0,
\end{equation}
then $L$ holds at every point of the smaller tube covered by both
initial strip charts.
\end{lemma}

\begin{proof}
Fix such a point $x$ and take one leaf of each strip through it,
using sheets whose preimages lie inside the buffered initial
blocks. Identify a smooth tube of the central torus with a disc
times that torus, and let $q$ be its smooth projection to the
torus. On a sufficiently large finite box in the universal plane
cover of the torus, the lifted projections of these two leaf
parametrizations are $C^1$ close to invertible affine maps.
The prescribed cylinder lengths may be chosen to include this
box and every lattice translate used below, with a fixed positive
buffer. These choices are made on the central model; the required
$C^1$ tolerance is then fixed on these finite buffered domains.
A small $C^1$ perturbation of an invertible affine map
is injective on a slightly smaller convex box, and its image
contains a further smaller box. We may therefore invert the
two projections and express the two leaf pieces as graphs over
one common convex box $Q$ in the torus cover. Choose this box to
contain a fundamental parallelogram for the two cycles and buffers
beyond every edge. The finite block lengths were chosen large
enough to accommodate these translates and buffers.

Each graph is periodic under its own primitive cycle wherever
both points and the required buffers belong to $Q$. To check
this exact assertion, shifting $w$ by $1$ leaves its strip map
unchanged. The lift of its torus projection changes by an
integral lattice vector; closeness to the original chart forces
that vector to be exactly the chosen cycle. Inversion and
uniqueness in the buffered box give the claimed graph equality.

Suppose the two holomorphic curve germs at $x$ were equal. The
set where the graph germs agree is open in $Q$. It is also
closed. At a limit of such points, both graph images have the
same limiting point. In small immersion charts at that point,
the two embedded holomorphic curve germs have accumulating
intersections, so the identity theorem makes them equal.
The use of the smooth projection causes no difficulty here:
it only supplies graph coordinates, while the identity theorem
is applied in holomorphic surface charts. Connectedness of $Q$
now makes the graphs agree throughout $Q$.

The common graph is
periodic under both edge identifications. It consequently
defines a compact smooth torus mapped into the tube whose local
image is a holomorphic curve. Let $m$ be the index of the
sublattice generated by the two cycles in the full period
lattice. Contraction of the disc component shows that the
represented integral homology class in $X_t$ is $m e$, with
$m>0$; in the application from Lemma~\ref{rat:safe}, $m=1$.
The pullback of $\sigma_t$ to a holomorphic curve is zero,
contradicting
\eqref{rat:destroy-fiber}. The two leaf germs are therefore
distinct.

They need not be transverse at $x$. Take a local foliation chart
for the second strip. Its transverse coordinate restricts to a
nonconstant holomorphic function on the first leaf, since the
two germs are distinct. Its derivative is nonzero at points
arbitrarily close to $x$. At such a point the first leaf meets
a member of the second strip family transversely. Apply
Remark~\ref{loc:alignment} at this intersection, assigning the
first strip, whose leaf contains $x$, to the role of $\sigma_2$
and the transverse partner to the role of $\sigma_1$.
Proposition~\ref{loc:criterion} then gives $L$ at $x$ by the
permitted displacement along the $\sigma_2$ leaf.
\end{proof}

\begin{proof}[Proof of Theorem~\ref{rat:open}]
By Lemma~\ref{rat:safe}, all central fibers except finitely many
have two independent safe cycle choices. Let $A\subset X_0$ be
the union of the excluded fibers and the singular fibers. This
is a proper analytic subset of dimension at most one.

We make a finite compact preparation before deforming. At every
point of $X_0\setminus A$, choose smaller regular torus patches
with two safe cycles and larger buffered patches. At a point of
$A$, choose a small holomorphic coordinate line disc through the
point whose boundary avoids $A$. Such a disc exists because a
generic line is not contained in $A$ and meets it discretely.
Two sufficiently small different tilts give transverse discs
with boundaries still outside $A$. Their centers may vary in a
smaller coordinate neighborhood. Compactness of $X_0$ permits
a finite choice of these center neighborhoods and regular
patches. Enlarge the compact regular data to include all the
disc boundaries. Every selected safe action then has a positive
common margin from its finite forbidden set.

Apply Lemma~\ref{rat:strips} to these finitely prepared data,
with blocks long enough for Lemma~\ref{rat:two-periods}.
On all sufficiently small deformations in $\mathcal D_v$, the
initial charts still cover fixed smaller tube patches, and the
coordinate-disc boundaries remain in those patches. If
\eqref{rat:destroy-fiber} holds, Lemma~\ref{rat:two-periods}
gives $L$ at every point of the required regular patches.
Corollary~\ref{loc:disc} gives $L$ in the reserved center
neighborhoods. Thus $X_t$ has $L$ everywhere.

Finally \eqref{rat:destroy-fiber} holds on a nonempty relatively
open subset arbitrarily close to $P_0$ in $\mathcal D_v$.
Indeed $(e,v)=0$, but the functional $(e,\cdot)$ does not vanish
identically on the positive vectors of $v^\perp$. Varying the
other normalized positive direction therefore breaks the
equation $(e,[\sigma_t])=0$. Local Torelli identifies our small
deformation neighborhood with a period neighborhood, giving
the desired nonempty relatively open set $\mathcal O_v$.
For any other marked K3 with the same period, global Torelli
after matching K\"ahler chambers identifies the underlying
surface, so the property has the stated period-theoretic scope.
\end{proof}
